\documentclass[twocolumn,10pt]{article}
\usepackage[margin=1in]{geometry}
\usepackage{amsmath,amssymb,amsthm}
\usepackage{hyperref}
\usepackage{booktabs}
\usepackage{listings}
\usepackage{xcolor}
\usepackage{microtype}
\usepackage{authblk}

\lstset{
  basicstyle=\ttfamily\footnotesize,
  breaklines=true,
  frame=single,
  backgroundcolor=\color{gray!8},
  captionpos=b
}

\newtheorem{proposition}{Proposition}

\title{\textbf{Detector-Based Oracles Are Empty Where\\
Quantum Error Correction Compilers\\
Must Be Validated}}

\author{Cleiton Augusto Correa Bezerra}
\affil{Independent Researcher\\
\texttt{augusto.cleiton@gmail.com}}
\date{September 2026}

\begin{document}
\maketitle

\begin{abstract}
A compiler pass over a quantum error correction circuit must preserve what the
circuit computes. The artefacts a practitioner reaches for when reasoning about
such a circuit---detector samples and the detector error model---are defined
relative to the circuit's own noiseless execution. We argue that this
self-reference makes them \emph{exactly} blind, rather than merely insensitive,
in the regime where a pass has to be checked: before noise is added, which is
also where the checking must happen if noise channels are not to survive the
gates they describe. Across 335 mutants that provably change what a
repetition-code memory circuit computes, detector sampling and detector error
model comparison detect none. On a rotated surface code memory experiment,
deleting every single-qubit gate leaves the detector error model byte-identical,
29 lines and 24 detectors either way, while 21 of 33 measurement outcomes change
by up to $0.50$. Two negative controls, each equivalent on every one of 200
trials with no oracle raising a false positive, establish that the measurement
is not an artefact of the mutation set. We give the measurement-record oracle
that closes the gap and is exact for noiseless stabiliser circuits, and we report
two negative results---one compiler test suite whose oracle detects every fault
we could construct, and one transformation library that came back clean---because
a method that only ever reports blindness is not a method anyone should trust.
All code and data are public.
\end{abstract}

\section{Introduction}

Quantum error correction has moved from demonstration to infrastructure, and
with it has come a layer of compilers: tools that lower a logical computation
onto a concrete circuit of resets, Clifford gates and measurements. Deltakit,
tqec and the transformation utilities inside Stim all perform rewrites that are
supposed to leave the computation unchanged.

The question this paper is about is narrow and practical. When such a rewrite is
applied, how does anyone check that it was correct?

The honest answer, in the projects we examined, is that the check varies widely
in strength, and that the artefact most readily available to a QEC practitioner
is the weakest of the options. That artefact is the detector: a parity of
measurement outcomes, declared by the circuit, which a decoder consumes. Its
aggregate form is the detector error model, which lists the error mechanisms a
noise model induces and which detectors each one flips.

Both are natural things to compare before and after a transformation. Both, we
argue, cannot see the transformation at all in the regime that matters.

\section{Why the blindness is structural}

A detector is not reported as a raw parity. It is reported as a deviation from
the value that parity takes when the circuit is executed without noise, and that
reference value is computed \emph{from the circuit itself}.

\begin{proposition}
Let $C$ be a stabiliser circuit with no noise channels, and let $C'$ be any
circuit obtained from $C$ by a transformation that leaves every detector
deterministic. Then every detector of $C$ and every detector of $C'$ reads zero,
each against its own reference, and no comparison of detector outcomes can
distinguish $C$ from $C'$.
\end{proposition}

The proposition is immediate once the self-reference is stated, and that is the
point: it is not a bound to be tightened by taking more shots. A detector-based
oracle applied to a noiseless circuit is not weak. It is empty.

The same argument reaches the detector error model. With no noise there are no
error mechanisms, so the model reduces to the detector declaration structure,
which a transformation that does not touch measurement records leaves alone.

One might object that this is a pathology of the noiseless case, and that real
circuits carry noise. The objection has the direction of the argument backwards.
A compiler simplification has to run \emph{before} noise is added, since a noise
channel attached to a gate that a later pass removes describes a gate that is no
longer in the circuit. The noiseless circuit is therefore not a toy. It is
exactly the object a pass is applied to, and exactly the object on which the
pass has to be shown correct.

\section{Method}

We follow the mutation-testing discipline of reporting, for each oracle, not how
many faults it caught but what fraction of the faults that mattered it could
have caught.

For every cell we report two numbers. The \emph{raw} rate is detections divided
by mutants applied. The \emph{power} is detections divided by mutants that
actually changed the circuit. The gap between them is the equivalent mutants:
faults that produce a circuit indistinguishable from the original, which no
correct oracle may flag and which counting as misses would understate every
oracle in the table.

Deciding which mutants are equivalent requires a reference oracle. On unitary
circuits the operator serves, and is exact. A QEC circuit has resets and
measurements, so it has no single unitary and that construction does not carry
over. What replaces it here is the measurement record. Our circuits are
noiseless and preserve the computational basis, so every measurement is
deterministic and a single shot settles equivalence exactly rather than
statistically. Where that property fails---a Hadamard before a measurement is
enough to break it---the record is random and the reference must be a
distribution over many shots, with a declared resolution. We flag this because
we got it wrong once: an early version of a related experiment compared single
shots on circuits containing Hadamards and produced 145 phantom divergences that
were nothing but different random draws.

\subsection{Negative controls}

A study that cannot distinguish its own controls is measuring nothing. We
therefore include two mutations that are equivalent by construction: removing a
gate immediately before a reset of the same qubit, and inserting a Pauli between
a measurement and the reset of that same ancilla. Both are dead by definition,
and any oracle that reports a difference on them is reporting a fault that is
not there.

An earlier version of this study had no such control and no Pauli pattern in its
circuits. Every qubit began in $|0\rangle$, nothing was noisy, and every
measurement read zero whatever the mutation did, so all four mutations were
reported as equivalent and the study said nothing at all.

\section{Results}

\subsection{Repetition-code memory}

Table~\ref{tab:main} reports 200 trials per mutation on repetition-code memory
at two widths and two round counts, with a random Pauli pattern applied after
the initial reset so that the syndrome carries information.

\begin{table}[t]
\centering
\small
\begin{tabular}{lrrrr}
\toprule
mutation & applied & equiv. & det. & DEM \\
\midrule
drop a dead gate \emph{(control)} & 200 & 200 & --- & --- \\
drop a CX & 200 & 34 & 0/166 & 0/166 \\
drop the control of a CX & 200 & 31 & 0/169 & 0/169 \\
insert a Pauli \emph{(control)} & 200 & 200 & --- & --- \\
\bottomrule
\end{tabular}
\caption{Detector sampling and detector error model comparison against mutations
of a compiled QEC circuit. Entries are detections over mutants that changed the
measurement record. Both controls are equivalent on every trial and no oracle
raises a false positive on them.}
\label{tab:main}
\end{table}

Over 335 mutants that provably change what the circuit computes, neither
detector-based oracle detects one.

\subsection{Rotated surface code}

On a distance-three rotated planar code memory experiment over three rounds, 17
qubits, compiled to a native gate set whose measurement is $M_X$ and whose reset
is $R_X$, we compare the compiled circuit against a mutant that removes every
single-qubit Clifford gate, all 65 of them.

The two detector error models are byte-for-byte identical: 29 lines, 24
detectors each. The measurement distributions differ in 21 of 33 positions, by
up to $0.50$.

A maximally destructive change to the unitary content of a surface code memory
circuit is therefore invisible to the artefact most often used to reason about
that circuit.

\subsection{With noise, on a matched pair}

The surface code comparison above is noiseless, and we do not extend it by
adding noise, because the mutant contains fewer gates and any gate-attached
noise model would then place a different number of channels on each side. That
difference, rather than the mutation, would drive the result.

On a single-qubit witness where the channel counts can be matched exactly---a
circuit of four instructions, with three depolarising channels on each side and
a deterministic $X$ inserted on one---the detector error models come out
identical to the last digit, \texttt{error(0.0066666666666666133) D0} on both,
and the detection rates over $2\times10^6$ shots agree to about $1.1\sigma$.
Blindness to a deterministic Pauli therefore survives the addition of noise
where the comparison can be made honestly.

\section{What we did not find}

Two negative results, reported because a method that only ever produces
blindness would be worth nothing.

\paragraph{tqec.} We examined the compilation test suite of \texttt{tqec}, a
topological QEC compiler. Its correctness assertions are three integers: the
code distance obtained from the shortest graphlike error of the noisy circuit,
the detector count, and the observable count. We expected this to be weak. It is
not. Across three classes of mutation it detected every mutant that changed the
circuit's behaviour, 40 of 40 in each, correctly ignored an equivalent mutation,
and raised no false positive on 25 identity insertions. The shortest graphlike
error is a structural property, not a counter, and breaking the circuit destroys
it. We report no defect.

\paragraph{Stim transformations.} Across 400 random Clifford circuits with
resets, measurements and repeat blocks, \texttt{flattened()}, string round-trip,
\texttt{without\_noise()} and concatenation all preserved the measurement
distribution. Deliberately broken transformations were caught 84 times out of
200 and 117 out of 156, so the clean result stands under partial and declared
sensitivity rather than under an unexamined one.

\section{What to use instead}

The oracle that works is the one the proposition does not reach: the measurement
record itself, compared between the two circuits rather than each against its
own reference. For a noiseless stabiliser circuit that preserves the
computational basis it is exact at one shot. Where the circuit is not
deterministic it becomes a distribution comparison with a declared resolution,
which is weaker but still strictly above anything detector-based.

This is cheap. It requires no decoder, no noise model, and no error budget. Its
only cost is that it must be done before noise is added, which is where the pass
belongs anyway.

\section{Reproducibility}

All code is public in the CleitonForge repository. The study of
Section~4.1 is \texttt{research/qec\_oracle\_power.py} and writes
\texttt{qec\_oracle\_power.json}. The Stim campaign of Section~5 is
\texttt{tools/fuzz\_stim\_transforms.py} and carries its own power measurement.
Versions: Stim 1.16.0, tqec 0.2.0.

\begin{thebibliography}{9}
\bibitem{gidney2021} C.~Gidney, \emph{Stim: a fast stabilizer circuit
simulator}, Quantum 5, 497 (2021).
\bibitem{cleitonforge} C.~A.~C.~Bezerra, \emph{CleitonForge}, \\
\url{https://github.com/cleitonaugusto/CleitonForge}
\bibitem{deltakit} Riverlane, \emph{Deltakit}, \\
\url{https://github.com/Deltakit/deltakit}
\bibitem{tqec} \emph{tqec: design automation for topological quantum error
correction}, \url{https://github.com/tqec/tqec}
\end{thebibliography}

\end{document}
